\section{Ergebnisse und Auswertung}\label{sec:experimente}

Dieses Kapitel beantwortet die zwei Forschungsfragen anhand der letzten
Datenmengenreihe. Die Modelle wurden ausschliesslich auf synthetischen
Bildern trainiert und auf dem aktuellen realen Validierungsset geprüft.
Zum Schluss wird die Erkennung in der App eingeordnet.

\subsection{Die gemeinsame Datenmengenreihe}\label{sec:saettigung}

Für die Reihe wurden sechs Datensätze mit 5'000, 10'000, 20'000, 50'000,
100'000 und 200'000 Bildern erzeugt. Bei jeder Bildanzahl wurden alle vier
Modelle der Ansätze A, B und C trainiert. Generatorstand, Seed,
Modellgrösse, Bildauflösung und Trainingsplan blieben gleich. Alle Modelle
wurden auf denselben 813 realen Frames ausgewertet. Davon zeigen 513
eine gelabelte Karte und 300 gelten als Negative.

Um zu prüfen, ob noch mehr Bilder weiterhelfen, wurde Variante~C
zusätzlich mit 500'000 Bildern trainiert. Generator, Seed,
Trainingsplan und Validierungsset sind dieselben. Der Lauf stammt von
einem späteren Commit, der nur den Speicherplatz der Trainingsmaschine
vergrösserte. Für A und B liegt bei dieser Bildanzahl kein Wert vor.

Die Datensätze sind keine ineinander enthaltenen Teilmengen. Die
Bildanzahl beeinflusst auch den Winkelplan und die Verteilung der
Kartenblätter (Abschnitt~\ref{sec:repro}). Innerhalb der Reihe lässt sich
dennoch untersuchen, wie sich die Leistung bei grösseren Datensätzen
entwickelt. Abbildung~\ref{fig:saettigung} zeigt die Ergebnisse.

\begin{figure}[H]
\centering
\includegraphics[width=0.92\textwidth]{saettigung.pdf}
\caption{Ergebnisse der sechs Datenmengen auf demselben realen
Validierungsset, für C zusätzlich mit 500'000 Bildern. Für C und B$_1$ wird \map{50} gezeigt, für A und B$_2$
die Top-1-Genauigkeit. Die Höhen dieser verschiedenen Kennzahlen sind
nicht direkt vergleichbar. B$_2$ erhält ideal entzerrte Ausschnitte.}
\label{fig:saettigung}
\end{figure}

Abbildung~\ref{fig:verlauf} zeigt den Trainingsverlauf von Variante~C.
Mit 5'000 Bildern verbessert sich das Modell während fast des gesamten
Trainings und erreicht erst gegen Ende sein bestes Ergebnis von 0.683.
Bei grösseren Datensätzen steigt der Wert schneller an. Ab 50'000
Bildern liegt die \map{50} bereits nach wenigen Epochen über 0.7 und
verändert sich danach nur noch wenig.

Bei den drei grössten Datensätzen wird das beste Ergebnis zwischen
Epoche 10 und 13 erreicht. Danach nimmt die Leistung wieder leicht ab.
Beim Lauf mit 200'000 Bildern endet sie bei 0.826 statt beim Höchstwert
von 0.908. Exportiert wird deshalb \code{best.pt} mit den Gewichten der
besten Epoche.

\begin{figure}[H]
\centering
\includegraphics[width=0.92\textwidth]{trainingsverlauf.pdf}
\caption{\map{50} der Variante~C auf dem realen Validierungsset je Epoche,
ein Feld je Datensatzgrösse. Der Punkt markiert die beste Epoche. Ab der
gestrichelten Linie ist die Mosaic-Augmentierung abgeschaltet.}
\label{fig:verlauf}
\end{figure}

\subsection{Vergleich der Erkennungsansätze (FF1)}\label{sec:vergleich}

Der einstufige Detektor C wird mit mehr Trainingsbildern bis 200'000
Bilder stetig besser und erreicht dort auf den realen Frames eine
\map{50} von 0.908. Bis hierhin passt der Verlauf zur Beobachtung von
Sun et al., dass die Leistung von Bildmodellen mit dem Logarithmus der
Datenmenge steigt \cite{sun2017revisiting}. Mit 500'000 Bildern liegt
die \map{50} bei 0.889. Mehr Bilder aus demselben Generator bringen
also keinen weiteren Gewinn. Der Abstand von 0.019 ist bei einem
einzelnen Seed nicht als Rückgang zu deuten
(Abschnitt~\ref{sec:setup}). Der Detektor gibt die Kartenklasse und ihre Position mit einem
Modell aus. Die reine Klassifikation A verbessert sich ebenfalls, liefert
aber keine Position der Karte.

Beim zweistufigen Ansatz B erkennt B$_1$ zuerst die Position der Karte.
Seine \map{50} verändert sich mit mehr Bildern nur wenig und liegt am Ende
der Reihe bei 0.842. B$_2$ klassifiziert einen ideal entzerrten
Kartenausschnitt fast immer richtig. Diese Ausschnitte stammen aus den von
Hand gesetzten Ecken. Im Betrieb muss B$_2$ dagegen mit den Vorhersagen von
B$_1$ arbeiten. Fehler beim Finden oder Entzerren der Karte wirken sich
dann auf die Klassifikation aus. Die letzte Reihe enthält keine
verlässliche Messung dieser gesamten Verarbeitungskette. Für B lässt
sich deshalb keine Trefferquote für die vollständige Erkennung angeben.

Für die mobile Nutzung spricht die vorliegende Evidenz am stärksten für C.
Ein weiterer Vorteil zeigt sich direkt bei der Bedienung. Die
achsenparallele Box umschliesst eine schräg liegende Karte zwar nicht
genau. Sie zeigt beim Zählen dennoch, welche Karte erkannt wurde und wo
sie sich im Bild befindet. Diese Rückmeldung fehlt bei A. Dadurch wirkt
eine Erkennung teilweise ungewohnt, weil ihr Bezug zum Kamerabild nicht
sichtbar ist. B benötigt zwei Modelle und einen zusätzlichen
Entzerrungsschritt. C übernimmt die Erkennung der Kartenklasse und ihrer
Position mit einem Modell und läuft in den eingereichten Apps auf
aktuellen Geräten in Echtzeit (Tabelle~\ref{tab:geraete}).
Ein neuer, gemeinsamer Latenzvergleich aller Varianten auf demselben Gerät
liegt für diese Reihe nicht vor. Der genaue Abstand beim Rechenaufwand
bleibt damit offen.

\subsection{Übertragung auf reale Aufnahmen (FF2)}\label{sec:transfer-real}

Das beste Ergebnis von C in der Datenmengenreihe stammt von einem Modell,
das keine realen Kartenhaufen als Trainingsbilder gesehen hat. Seine
\map{50} von 0.908 auf den realen Aufnahmen zeigt, dass die synthetischen
Bilder für die Erkennung auf diesem Validierungsset geeignet sind. Die
Kennzahl berücksichtigt die richtige Kartenklasse und die Überlappung
der erkannten Position mit dem Label. Sie ist nicht gleichbedeutend mit
dem Anteil richtig gezählter Karten in einer Zählsession.

Im Projektverlauf wurden Generator und reale Prüfbasis erweitert.
Insbesondere flache Kamerawinkel erhielten mehr Gewicht. Der Generator
deckt nun Neigungen bis 60\textdegree{} ab. Hinzu kamen Aufnahmen unter
anderen Lichtverhältnissen, beide Kartenblätter und schwierigere
Negativbeispiele. Die reale Prüfbasis enthält ebenfalls mehr flache
Perspektiven und Situationen ohne klar lesbare Karte. Diese Änderungen
verbessern die Abdeckung der vorgesehenen Einsatzbedingungen. Ihr
einzelner Beitrag zur Modellleistung wurde nicht getrennt gemessen.

Die Aussage zum Transfer gilt zunächst nur für die vorhandenen
Aufnahmen. Viele Frames stammen aus denselben Videos und ähneln sich.
Ausserdem wurde die beste Trainingsepoche anhand dieses
Validierungssets gewählt. Ein unabhängiges Testset und Wiederholungen
mit weiteren Trainings-Seeds fehlen. Besonders bei Karten im Wurf oder
an der Grenze zwischen lesbar und unlesbar bleibt auch die Bewertung
schwierig (Abschnitt~\ref{sec:grenzfaelle}).

\subsection{Betrieb in der App}\label{sec:betriebspunkt}

Die eingereichten Apps verwenden Variante~C. Eine Karte wird erst
gezählt, wenn die Erkennung die eingestellte Konfidenz erreicht und über
mehrere Frames stabil bleibt (Abschnitt~\ref{sec:zaehlung}). Eine
aufgezeichnete Zählsession wurde im Dataset Tool mit verschiedenen
Einstellungen erneut ausgewertet. Dabei zeigte sich, dass einzelne
Fehler auch mit einer hohen Konfidenz auftreten können
(Abbildung~\ref{fig:frames-nacheinander}). Eine höhere Schwelle allein
löst dieses Problem daher nicht. Die Stabilitätsregel hilft, solche
kurzen Ausreisser abzufangen.

\begin{figure}[H]
\centering
\includegraphics[width=0.48\textwidth]{detection-difficulties-wrong.png}\hfill
\includegraphics[width=0.48\textwidth]{detection-difficulties-correct.png}
\caption{Zwei direkt aufeinanderfolgende Frames derselben Aufnahme. Links
erkennt das Modell die oberste Karte mit 74\,\% fälschlich als Schellen~9.
Diese Erkennung wird nicht gezählt. Im nächsten Frame erkennt es die
Schellen~10 mit 87\,\% korrekt und die App zählt die Karte. Die gelbe
Markierung zeigt die gezählte Erkennung. Der Kartenhaufen und die Position der
Kamera haben sich zwischen den beiden Frames kaum verändert.}
\label{fig:frames-nacheinander}
\end{figure}

Diese einzelne Session reicht nicht für eine verlässliche Erfolgsquote.
Sie zeigt jedoch, wie sich die Einstellungen der App anhand einer
vollständigen Aufnahme prüfen lassen. Für aussagekräftige Ergebnisse
braucht es weitere Sessions mit unterschiedlichen Geräten und
Spielsituationen. Dabei sollten ganze Zählvorgänge statt nur einzelne
Frames ausgewertet werden.
